%% =========================================================================
\section{The Min/Max Dichotomy: No Congruence Invariant Blocks a Prime}
\label{sec:min-max-dichotomy}
%% =========================================================================

Every theorem asserting that a prime is \emph{missing} from a
Euclid--Mullin sequence has, historically, been a theorem about the
$\mathrm{maxFac}$ rule.  This section makes the difference between the
two rules into a formal statement.  On the $\mathrm{maxFac}$ side we
machine-verify a complete omission certificate---the prime~$5$ never
occurs in the second sequence---and package it as an instance of an
abstract obstruction class.  On the $\minFac$ side we prove that the
very same obstruction class is \emph{empty}: for a missing prime~$q$
and \emph{every} modulus rich enough to express the capture condition
for~$q$, no propagating congruence invariant can block~$q$.

Two further layers sharpen this.  Abstracting the argument away from
residues isolates the single semantic condition that drives it
(\S\ref{sec:obstruction-calculus}), and lifting it from invariants to
\emph{proofs} turns the no-go into a statement about a proof system:
within the congruence-induction fragment, provability of avoidance and
membership in the sequence are the same predicate
(\S\ref{sec:proof-fragment}).

\begin{roadmap}
\begin{enumerate}[nosep,leftmargin=*]
\item[\S\ref{sec:pop-to-point}] The population-to-point gap --- inverse theorems and obstruction-killing.
\item[\S\ref{sec:two-sequences}] The two sequences --- open congruence condition vs.\ thin multiplicative condition.
\item[\S\ref{sec:cvdp-control}] The control case --- the Cox--van der Poorten obstruction for $\mathrm{maxFac}$, formalized.
\item[\S\ref{sec:no-invariant}] The No-Invariant Theorem --- eviction, free-state fullness, CRT-reach.
\item[\S\ref{sec:formal-dichotomy}] The formal dichotomy --- the obstruction class is inhabited for max, empty for min.
\item[\S\ref{sec:ic-min}] $\mathrm{IC}_{\min}$ and its exact position in the reduction network.
\item[\S\ref{sec:obstruction-calculus}] The obstruction calculus --- enrichments, certificates, the Emptiness Theorem; max provably not killable.
\item[\S\ref{sec:proof-fragment}] The proof-theoretic dichotomy --- provability of avoidance decides appearance; every guard direction closed.
\item[\S\ref{sec:reciprocity-frontier}] The reciprocity frontier --- no invariant at the growing symbol modulus blocks a prime.
\end{enumerate}
\end{roadmap}

%% =========================================================================
\subsection{The Population-to-Point Gap}
\label{sec:pop-to-point}
%% =========================================================================

The recurring difficulty of this paper is that MC asks for a
\emph{generic} property (equidistribution, hitting) of one
\emph{specific} object (the orbit of~$2$).  Sections~\ref{sec:walk}
through~\ref{sec:hard} attack the positive half of that problem:
find a hypothesis under which the specific orbit inherits the generic
behaviour.

There is a complementary half, standard in additive combinatorics and
in the modern treatment of pseudorandomness, that this section
supplies.  The technique has two moves:

\begin{enumerate}[nosep]
\item An \textbf{inverse theorem}: failure of the generic behaviour
  forces the object to correlate with a \emph{classified} obstruction.
\item \textbf{Obstruction-killing}: the specific object provably
  correlates with none of the classified obstructions.
\end{enumerate}

\noindent
Together these give the generic behaviour for the specific object.  We
are not aware of any inverse theorem for either Euclid--Mullin sequence
in the literature; the search for one, in both directions, came back
empty.  What this section delivers is move~2, in full, for one
obstruction class: the class of \emph{propagating congruence
invariants}.  This is exactly the class into which the classical
$\mathrm{maxFac}$ omission proof of Cox and van der
Poorten~\cite{CoxVdP1968} falls.  The result is a no-go theorem, and
its scientific content is entirely negative: it says where a disproof
of MC cannot come from.

%% =========================================================================
\subsection{The Two Sequences and the Structural Asymmetry}
\label{sec:two-sequences}
%% =========================================================================

The second Euclid--Mullin sequence (OEIS A000946) replaces the smallest
prime factor by the largest:
\[
  Q_0 = 2, \qquad
  r_n = \mathrm{maxFac}(Q_{n-1}+1), \qquad
  Q_n = Q_{n-1} \cdot r_n ,
\]
formalized as
\lean{EM/Obstruction/MaxVariant.lean}{mseq} and
\lean{EM/Obstruction/MaxVariant.lean}{mprod}, built from
\lean{EM/Obstruction/MaxVariant.lean}{maxFac},
the supremum of \code{n.primeFactors}.  The basic structure theory
mirrors the $\minFac$ side: every term is prime
(\lean{EM/Obstruction/MaxVariant.lean}{mseq\_prime}), no prime repeats
(\lean{EM/Obstruction/MaxVariant.lean}{mseq\_injective}), and every term
divides every later accumulator
(\lean{EM/Obstruction/MaxVariant.lean}{mseq\_dvd\_mprod}).
Cox and van der Poorten~\cite{CoxVdP1968} showed that the primes
$\{5, 11, 13, 17, 19, 23, 29, 31, 37, 41, 47\}$ (all primes below~$53$
other than the terms $2,3,7,43$) are missing from this sequence, and
conjectured that infinitely many primes are.  Booker~\cite{Booker2012} showed that infinitely many primes are
omitted; Pollack and Treviño~\cite{PollackTrevino2014} gave an
elementary proof of Booker's theorem.

\paragraph{Where the two rules part company.}
Fix a prime~$q$ and a Euclid number~$N$.  Ask: what does it take for
the rule to \emph{select}~$q$ at this step?

\smallskip
\begin{center}
\begin{tabular}{@{}l@{\quad}l@{}}
\toprule
\textbf{Rule} & \textbf{Appearance condition} $\sigma(N) = q$ \\
\midrule
$\minFac$
  & $q \mid N$ and no prime $< q$ divides~$N$ \\
  & \emph{an open congruence condition} \\[0.4em]
$\mathrm{maxFac}$
  & $q \mid N$ and $N$ is $q$-smooth \\
  & \emph{a thin multiplicative condition} \\
\bottomrule
\end{tabular}
\end{center}

\smallskip\noindent
The first condition is a nonempty union of residue classes modulo
$q \cdot \prod_{p < q,\ p \text{ odd}} p$---the modulus called
\lean{EM/Obstruction/NoInvariant.lean}{forcingModulus} in the
formalization.  It is a positive-density congruence condition that meets
\emph{every} residue class mod~$4$, and it says nothing at all about the
large prime factors of~$N$: the cofactor $N/q$ is arithmetically
unconstrained.

The second condition is not a congruence condition at any modulus:
every residue class contains integers with arbitrarily large prime
factors, so $q$-smoothness cannot be detected modulo anything.  But
intersected with congruence data it collapses violently.  For $q = 5$,
smoothness plus ``$N$ odd'' plus ``$3 \nmid N$'' forces
$N = 5^c$ (\lean{EM/Obstruction/MaxVariant.lean}{eq\_pow\_five\_of\_maxFac\_five}),
a single geometric progression, all of whose members lie in the single
class $1 \pmod 4$
(\lean{EM/Obstruction/MaxVariant.lean}{five\_pow\_mod\_four}).
A congruence invariant now has something to bite on.

This is the whole asymmetry: \emph{congruence invariants can kill thin
multiplicative conditions, and have nothing to bite on against open
congruence conditions.}  The formalization witnesses it on a single
number.

\begin{theorem}[The asymmetry witness ---
\lean{EM/Obstruction/MaxVariant.lean}{min\_side\_no\_smoothness\_forcing},
\lean{EM/Obstruction/MaxVariant.lean}{cvdp\_selection\_rule\_asymmetry}]
\label{thm:asymmetry-witness}
There is a natural number~$N$ satisfying \emph{every} congruence
hypothesis used in the Cox--van der Poorten argument---$N \geq 3$,
$N$ odd, $3 \nmid N$, $N \equiv 3 \pmod 4$---with
$\minFac(N) = 5$ and $\mathrm{maxFac}(N) \neq 5$; moreover $N$ is not a
power of~$5$.  The witness is $N = 35$.  (Since
\S\ref{sec:formal-dichotomy} the single-number witness has been superseded by
a statement about the entire class of arguments; it is kept here because it
exhibits the mechanism on one integer.)
\end{theorem}

\noindent
The same number is a forcing state for~$5$ on the min side and is
blocked on the max side.  Note that
\lean{EM/Obstruction/MaxVariant.lean}{maxFac\_thirtyfive\_ne\_five}
is derived abstractly from the mod-$4$ collision
(\lean{EM/Obstruction/MaxVariant.lean}{maxFac\_ne\_five\_of\_mod\_four}),
not by factoring~$35$.

%% =========================================================================
\subsection{The Control Case: the Cox--van der Poorten Obstruction, Formalized}
\label{sec:cvdp-control}
%% =========================================================================

Before proving that an object does not exist, it is worth building one.
The formalization carries a self-contained construction of the second
sequence and a machine-checked omission theorem for it.

\begin{keyresult}
\begin{theorem}[\lean{EM/Obstruction/MaxVariant.lean}{five\_not\_mem\_mseq}]
\label{thm:five-not-mem}
The prime $5$ never occurs in the second Euclid--Mullin sequence:
$\forall n,\ \mathrm{mseq}(n) \neq 5$.  Unconditional, computation-free,
zero \code{sorry}.
\end{theorem}
\end{keyresult}

\noindent The proof is three moves.

\begin{enumerate}[nosep]
\item \emph{Congruence bookkeeping.}  The prime~$3$ occurs at step~$1$
  (\lean{EM/Obstruction/MaxVariant.lean}{mseq\_one}), so
  $3 \mid Q_n$ for $n \geq 1$
  (\lean{EM/Obstruction/MaxVariant.lean}{three\_dvd\_mprod}) and hence
  $3 \nmid Q_n + 1$
  (\lean{EM/Obstruction/MaxVariant.lean}{three\_not\_dvd\_mprod\_succ}).
  Every multiplier past step~$0$ is odd
  (\lean{EM/Obstruction/MaxVariant.lean}{mseq\_succ\_odd}), whence
  $Q_n \equiv 2 \pmod 4$
  (\lean{EM/Obstruction/MaxVariant.lean}{mprod\_mod\_four}) and, combining,
  $Q_n \equiv 6 \pmod{12}$
  (\lean{EM/Obstruction/MaxVariant.lean}{mprod\_mod\_twelve}).
  So the Euclid number satisfies $N \equiv 3 \pmod 4$ and $3 \nmid N$.
\item \emph{Smoothness forcing.}  $\mathrm{maxFac}(N) = 5$ bounds
  \emph{all} prime factors of~$N$ by~$5$; oddness kills~$2$ and step~1
  kills~$3$, so $N = 5^c$
  (\lean{EM/Obstruction/MaxVariant.lean}{eq\_pow\_five\_of\_maxFac\_five}).
\item \emph{The mod-$4$ collision.}  $5^c \equiv 1 \pmod 4$ for
  every~$c$, contradicting $N \equiv 3 \pmod 4$
  (\lean{EM/Obstruction/MaxVariant.lean}{maxFac\_ne\_five\_of\_mod\_four}).
\end{enumerate}

\paragraph{The obstruction, packaged.}
The three moves are then repackaged abstractly, so that the omission
genuinely \emph{flows from} an obstruction rather than merely
coinciding with one.  Define the max-side transition relation on
$\ZZ/m\ZZ$
(\lean{EM/Obstruction/MaxVariant.lean}{MaxR}):
$r \to r'$ iff there is an odd $N \geq 3$ with
$N \equiv r+1$ and $r' = r \cdot \mathrm{maxFac}(N)$.
A set is \emph{propagating}
(\lean{EM/Obstruction/MaxVariant.lean}{MaxPropagating}) if closed under
this relation; a state is \emph{forcing} for~$q$
(\lean{EM/Obstruction/MaxVariant.lean}{MaxForcingState}) if some legal
candidate in its class has largest prime factor~$q$; and $S$
\emph{blocks}~$q$
(\lean{EM/Obstruction/MaxVariant.lean}{MaxBlocks}) if it contains no
forcing state.  A
\lean{EM/Obstruction/MaxVariant.lean}{MaxCvdPObstruction} is a propagating
set that blocks~$q$ and contains the orbit tail.

\begin{keyresult}
\begin{theorem}[The obstruction class is inhabited ---
\lean{EM/Obstruction/MaxVariant.lean}{max\_cvdp\_obstruction\_five}]
\label{thm:max-obstruction-inhabited}
The singleton $\{6\} \subseteq \ZZ/12\ZZ$ is a Cox--van der Poorten
obstruction to the prime~$5$ for the $\mathrm{maxFac}$ sequence.
\end{theorem}
\end{keyresult}

\noindent
Propagation is
\lean{EM/Obstruction/MaxVariant.lean}{maxPropagating\_six}
(an odd $N \equiv 7 \bmod 12$ has odd largest prime factor~$p$, and
$6p \equiv 6 \bmod 12$); blocking is
\lean{EM/Obstruction/MaxVariant.lean}{maxBlocks\_six\_five}
(exactly moves~2 and~3 above); tail containment is
\lean{EM/Obstruction/MaxVariant.lean}{mprod\_mem\_six}.  The real sequence
step is an instance of the abstract relation
(\lean{EM/Obstruction/MaxVariant.lean}{mprod\_step\_maxR}), and the
omission is then re-derived abstractly: any prime blocked by a
propagating set containing the orbit tail is missing
(\lean{EM/Obstruction/MaxVariant.lean}{missing\_of\_obstruction}), giving
\lean{EM/Obstruction/MaxVariant.lean}{five\_not\_mem\_mseq'}.

\begin{remark}[Why the modulus must be $12$, not $4$]
\label{rem:mod-twelve}
The mod-$4$ collision supplies the \emph{contradiction}, but the mod-$3$
information is essential to \emph{propagation and blocking}.  At $m = 4$
the singleton $\{2\}$ does not block~$5$: the number $N = 15$ is odd,
$\equiv 3 \pmod 4$, and has $\mathrm{maxFac}(15) = 5$, so the state
$2 \in \ZZ/4\ZZ$ admits a legal candidate capturing~$5$.  Only after
$3 \nmid N$ is imposed---which costs the factor~$3$ in the
modulus---does smoothness collapse to $\{5^c\}$.  \emph{(The failure at
$m=4$ is stated here for orientation only; it is not part of the Lean
development, which works throughout at $m = 12$.)}
\end{remark}

\begin{remark}[Reach of this reconstruction]
\label{rem:reach-mod-four}
The mod-$4$ collision closes $q = 5$ and only $q = 5$.  For $q = 11$ the
forced smooth numbers are $5^a 7^b 11^c$, and since
$7 \equiv 11 \equiv 3 \pmod 4$ these do meet the class $3 \pmod 4$;
mod~$4$ alone gives no contradiction.  The remaining Cox--van der
Poorten primes need further quadratic-residue input, and no claim about
them is made here.
\end{remark}

%% =========================================================================
\subsection{The No-Invariant Theorem}
\label{sec:no-invariant}
%% =========================================================================

We now mirror the entire apparatus on the $\minFac$ side and prove it
vacuous.

\begin{definition}[The min-side transition system]
\label{def:min-transition}
Fix a modulus $m \neq 0$.  The \defname{transition relation}
on $\ZZ/m\ZZ$
(\lean{EM/Obstruction/NoInvariant.lean}{Transition}) is
\[
  r \to r'
  \iff
  \exists N \in \NN,\quad
  N \text{ odd},\ N \geq 3,\
  N \equiv r + 1 \ (\mathrm{mod}\ m),\
  r' = r \cdot \minFac(N).
\]
A set $S \subseteq \ZZ/m\ZZ$ is \defname{propagating}
(\lean{EM/Obstruction/NoInvariant.lean}{Propagating}) if it is closed
under~$\to$; $r$ is a \defname{forcing state} for~$q$
(\lean{EM/Obstruction/NoInvariant.lean}{ForcingState}) if \emph{every}
candidate~$N$ in the class $r+1$ has $q \mid N$ and is divisible by no
odd prime below~$q$; $S$ \defname{blocks}~$q$
(\lean{EM/Obstruction/NoInvariant.lean}{Blocks}) if it contains no forcing
state; and $S$ \defname{contains the tail}
(\lean{EM/Obstruction/NoInvariant.lean}{ContainsTail}) if
$\Prod{2}(n) \bmod m \in S$ for all large~$n$.
A \defname{CvdP obstruction}
(\lean{EM/Obstruction/NoInvariant.lean}{CvdPObstruction}) is a set that is
propagating, blocks~$q$, and contains the tail.
The modulus is \defname{rich enough}
(\lean{EM/Obstruction/NoInvariant.lean}{RichEnough}) for~$q$ if
$q \mid m$ and every odd prime $< q$ divides~$m$; the minimal such
modulus is
$M(q) = q \cdot \prod_{p < q,\ p \text{ odd}} p$
(\lean{EM/Obstruction/NoInvariant.lean}{forcingModulus},
\lean{EM/Obstruction/NoInvariant.lean}{richEnough\_of\_forcingModulus\_dvd}).
\end{definition}

\begin{remark}[Deliberate over-approximation, and what it costs]
\label{rem:over-approximation}
In \code{Transition}, $N$ ranges over \emph{all} natural numbers in the
residue class $r+1$---not only those actually realizable as
Euclid--Mullin candidates.  The relation is therefore strictly larger
than the true orbit dynamics.  We should be precise about the direction
this cuts.  If $\to_{\mathrm{true}} \subseteq \to_{\mathrm{big}}$, then a set that
propagates under $\to_{\mathrm{big}}$ propagates under
$\to_{\mathrm{true}}$, so ``no big-propagating blocking set exists'' is
\emph{implied by}, hence \emph{weaker than}, ``no true-propagating
blocking set exists''.  (For the true dynamics the latter question is in
fact vacuous: the tail residues of the actual orbit form a propagating
set, which blocks~$q$ iff $q$ is missing.)  The honest reading of the
No-Invariant Theorem is therefore about \emph{proofs}, not dynamics: it
excludes any omission argument whose only knowledge of the candidate~$N$
is its residue class modulo~$m$---because, on the min side, capture
$\minFac N = q$ is itself a residue-class condition on~$N$
(Theorem~\ref{thm:min-congruence-determined}) and Dirichlet supplies
candidates in every class.  That is a one-paragraph observation with a
machine-checked proof, and it is all we claim.

The two definitions of ``blocks'' are set up in the same spirit: on the
min side \code{ForcingState} is universally quantified over candidates,
on the max side \code{MaxForcingState} is existentially quantified, so
that each side is set up against the conclusion it reaches.
\end{remark}

\noindent
That \code{Blocks} is the right notion is not a matter of taste; it is a
lemma.  At a forcing state the capture really happens:

\begin{lemma}[Semantics of forcing ---
\lean{EM/Obstruction/NoInvariant.lean}{forcingState\_captures}]
\label{lem:forcing-captures}
If $q$ is prime and $\Prod{2}(n) \bmod m$ is a forcing state for~$q$,
then $\seq{2}(n{+}1) = q$.
\end{lemma}

\noindent
So a set containing a forcing state cannot certify that $q$ is missing;
a genuine certificate must block, and the definition captures exactly
that.

\paragraph{Move 1: eviction.}
The engine is the finiteness of the hitting set for a missing prime.

\begin{lemma}[Finite Hitting ---
\lean{EM/Obstruction/NoInvariant.lean}{hittingSet\_finite},
alias \lean{EM/Obstruction/NoInvariant.lean}{hittingSet\_finite}]
\label{lem:finite-hitting}
For $q$ a missing prime, $\mathrm{HittingSet}(q) = \{n : q \mid
\Prod{2}(n)+1\}$ is finite, with
$|\mathrm{HittingSet}(q)| \leq q$
(\lean{EM/Obstruction/NoInvariant.lean}{hittingSet\_ncard\_le}).
\end{lemma}

\begin{proof}[Proof sketch]
At a hitting step the guardian bound
(\lean{EM/Population/HittingSetStructure.lean}{hitting\_step\_guardian})
gives $\seq{2}(n{+}1) < q$, since $\minFac(\Prod{2}(n)+1) \le q$ and
equality would mean $q$ appears.  Because no prime repeats
(\lean{EM/Core/Defs.lean}{seq\_injective}), the map
$n \mapsto \seq{2}(n{+}1)$ is injective on the hitting set, which
therefore injects into the primes below~$q$.
\end{proof}

\noindent
This is the first place the rule matters: the guardian bound
$\minFac(\Prod{2}(n)+1) < q$ is exactly what fails for
$\mathrm{maxFac}$, where the chosen factor is $\geq q$ and the hitting
set may well be infinite.

Finite hitting feeds an eviction statement that is
\textbf{unconditional}---no missing-prime assumption:

\begin{lemma}[Eviction ---
\lean{EM/Obstruction/NoInvariant.lean}{exists\_tail\_coprime}]
\label{lem:eviction}
For every modulus $m \neq 0$ there is $N_0$ with
$\gcd(\Prod{2}(n)+1, m) = 1$ for all $n \geq N_0$.
\end{lemma}

\noindent
The proof is the \textbf{A/B split}, applied to each prime $\ell \mid m$
(\lean{EM/Obstruction/NoInvariant.lean}{eventually\_not\_dvd\_succ}):
either $\ell$ occurs in the sequence, in which case $\ell \mid
\Prod{2}(n)$ from that point on and so $\ell \nmid \Prod{2}(n)+1$; or
$\ell$ is missing, in which case its hitting set is finite by
Lemma~\ref{lem:finite-hitting}.  Taking the maximum over the finitely
many prime factors of~$m$ gives $N_0$.  In particular the orbit is
eventually \emph{free} forever: $\Prod{2}(n)+1$ is a unit mod~$m$.

\paragraph{Move 2: free-state fullness.}
From a free state, every unit multiple is reachable in a single step.

\begin{lemma}[Free-state fullness ---
\lean{EM/Obstruction/NoInvariant.lean}{free\_transition}]
\label{lem:free-transition}
Let $m \neq 0$, $r \in \ZZ/m\ZZ$ with $r+1$ a unit.  Then for
\emph{every} unit $s \in (\ZZ/m\ZZ)^\times$ there is a transition
$r \to r\cdot s$.
\end{lemma}

\begin{proof}[Proof sketch]
Apply Dirichlet twice.  Pick a prime $\pi > 2$ with $\pi \equiv s$, then
a prime $M > \pi$ with $M \equiv (r+1)s^{-1}$---a unit class, since both
$r+1$ and $s$ are units.  Set $N = \pi M$.  Then $N$ is odd,
$N \equiv r+1$, and $\minFac(N) = \pi \equiv s$
(\lean{EM/Obstruction/NoInvariant.lean}{minFac\_mul\_of\_lt}).
\end{proof}

\noindent
The Mathlib input is
\code{Nat.forall\_exists\_prime\_gt\_and\_eq\_mod}
(\code{Mathlib/NumberTheory/LSeries/PrimesInAP.lean}).
This lemma is \textbf{rule-symmetric}, and it is \emph{not} where the
two rules part company.  For $\mathrm{maxFac}$ the very same
$N = \pi M$ would give $\mathrm{maxFac}(N) = M$, whose class is
$(r+1)s^{-1}$; but one is free to pick the two primes in the other
order---first a prime $\pi \equiv (r+1)s^{-1}$, then a prime
$M > \pi$ with $M \equiv s$, so that again $N = \pi M \equiv r+1$
while now $\mathrm{maxFac}(N) = M \equiv s$.  Free-state fullness
therefore holds under \emph{both} rules.  The genuine break points are
the guardian bound (Lemma~\ref{lem:finite-hitting}) and, decisively,
the capture condition of \S\ref{sec:two-sequences}---a congruence
condition for $\minFac$, satisfiable modulo
$\mathrm{forcingModulus}(q)$, versus a smoothness condition for
$\mathrm{maxFac}$, which no fixed modulus can express, so max-side
blocking degenerates and obstructions exist trivially.  That is where
Booker-type omissions live; see
\S\ref{sec:congruence-receptacle}.

\begin{remark}[Consistency note]
\label{rem:fullness-symmetric}
For a time our own reading---and the Lean docstrings---located the
min/max break point at this lemma: under $\minFac$ the pair
$N = \pi M$ with $\pi < M$ gives multiplier~$\pi$, whose class is
free, while under $\mathrm{maxFac}$ it gives~$M$, whose class is
forced to be $(r{+}1)s^{-1}$.  That reading is \textbf{wrong}, and
the docstrings have been corrected accordingly: as shown above, one
merely picks the two primes in the other order.  The rows of the
summary table in \S\ref{sec:formal-dichotomy} recording the
multiplier of $\pi M$ and the reachable set should be read with this
correction in mind: they describe the \emph{particular} candidate
used in the min-side proof, not an obstruction on the max side.  The
two rules are separated by the capture condition and by the guardian
bound (Lemma~\ref{lem:finite-hitting}), not by fullness.
\end{remark}

\paragraph{Move 3: CRT-reach.}
Combining the moves gives the theorem.

\begin{keyresult}
\begin{theorem}[No-Invariant Theorem ---
\lean{EM/Obstruction/NoInvariant.lean}{no\_cvdp\_obstruction}]
\label{thm:no-invariant}
Let $q$ be a prime that never occurs in the (first, $\minFac$)
Euclid--Mullin sequence.  Then for \emph{every} modulus $m \neq 0$ with
$\mathrm{RichEnough}(q,m)$, \textbf{no} set $S \subseteq \ZZ/m\ZZ$ is a
CvdP obstruction for~$q$ modulo~$m$.
\end{theorem}
\end{keyresult}

\begin{proof}[Proof sketch]
By Lemma~\ref{lem:eviction} the tail is free forever, and by
\code{ContainsTail} some free tail state $r = \Prod{2}(n) \bmod m$ lies
in~$S$.  By Lemma~\ref{lem:free-transition} and propagation, $S$
contains $r \cdot s$ for every unit~$s$.  Choose by CRT a unit~$s$ with
$s \equiv -\Prod{2}(n)^{-1} \pmod q$ and $s \equiv 1$ modulo every other
prime factor of~$m$.  Then $r s + 1 \equiv 0 \pmod q$, while for every
odd prime $p < q$ we get $rs + 1 \equiv r + 1 \not\equiv 0 \pmod p$
\emph{by freeness}---so the A/B split is not even needed at this stage.
Thus $r \cdot s \in S$ is a forcing state, contradicting
\code{Blocks}.
\end{proof}

\noindent
No parity hypothesis on~$m$ is needed.  The frozen $2$-adic data of the
first sequence ($\Prod{2}(n) \equiv 2 \bmod 4$, candidate
$\equiv 3 \bmod 4$) blocks nothing, because ``$q \mid N$ and
$N \equiv 3 \pmod 4$'' is satisfiable for every odd~$q$, and the
Dirichlet construction prescribes the candidate class modulo the
\emph{whole} modulus, absorbing the mod-$4$ constraint automatically.
The two-part specialisation is
\lean{EM/Obstruction/NoInvariant.lean}{no\_cvdp\_obstruction\_two\_part}:
for $m = 2^a m_0$ with $m_0 \neq 0$, still no obstruction.

Contrast with max, where the appearance condition for~$5$ is $N = 5^c$,
which forces $N \equiv 1 \pmod 4$ and is killed by the frozen
$3 \bmod 4$.  \emph{That} is the dichotomy: a congruence obstruction can
only bite when the appearance condition is not itself a full-freedom
congruence condition, and for $\minFac$ it always is.

\paragraph{Lifting, with an honest caveat.}
An invariant modulo~$m$ pulls back along $\ZZ/m'\ZZ \to \ZZ/m\ZZ$
whenever $m \mid m'$
(\lean{EM/Obstruction/NoInvariant.lean}{liftSet}).  Propagation lifts
(\lean{EM/Obstruction/NoInvariant.lean}{propagating\_lift}) and tail
containment lifts
(\lean{EM/Obstruction/NoInvariant.lean}{containsTail\_lift}), but the
forcing form of \code{Blocks} does \emph{not}: \code{ForcingState} at the
finer modulus is a weaker condition on the projected state, since the
class being quantified over is smaller.  What does lift is the stronger
\emph{death-avoiding} form
(\lean{EM/Obstruction/NoInvariant.lean}{BlocksDeath}), in which $S$ avoids
$r+1 \equiv 0 \bmod q$ altogether; it implies \code{Blocks}
(\lean{EM/Obstruction/NoInvariant.lean}{blocks\_of\_blocksDeath}) and it
lifts (\lean{EM/Obstruction/NoInvariant.lean}{blocksDeath\_lift}), giving
\lean{EM/Obstruction/NoInvariant.lean}{cvdpObstruction\_lift}.
This is why Theorem~\ref{thm:no-invariant} takes
$\mathrm{RichEnough}(q,m)$ as a hypothesis rather than deriving it by
lifting from a coarser modulus.

%% =========================================================================
\subsection{The Formal Dichotomy}
\label{sec:formal-dichotomy}
%% =========================================================================

Putting the two halves side by side gives the centerpiece.

\begin{keyresult}
\textbf{The min/max dichotomy}
(\lean{EM/Obstruction/RuleTransition.lean}{cvdp\_dichotomy}).
Let $\Phi$ be a selection rule and write
\[
  R_\Phi^{(m)}(r, r') \;:\Longleftrightarrow\;
  \exists N \text{ odd},\ N \geq 3,\ N \equiv r+1 \!\!\pmod m,\
  r' = r \cdot \Phi(N),
\]
with \emph{propagating}, \emph{forcing} and \emph{blocking} defined from it and
an accumulator supplying the orbit tail.  For this one system:
\begin{itemize}[nosep]
\item the class is \textbf{inhabited} at $\Phi = \mathrm{maxFac}$:
  $\{6\} \subseteq \ZZ/12\ZZ$ blocks~$5$
  (\lean{EM/Obstruction/MaxVariant.lean}{max\_cvdp\_obstruction\_five}), and
  the omission of~$5$ follows from it
  (\lean{EM/Obstruction/MaxVariant.lean}{five\_not\_mem\_mseq'});
\item the class is \textbf{empty} at $\Phi = \minFac$, for every missing
  prime~$q$ and \emph{every} modulus
  (\lean{EM/Obstruction/RuleTransition.lean}{no\_min\_rule\_obstruction}).
\end{itemize}
The two sides are instances of the same definition, not analogues of it:
$\mathrm{MaxR}$ and $\mathrm{Transition}$ are $R_\Phi$ at
$\Phi = \mathrm{maxFac}$ and $\Phi = \minFac$ by \texttt{Iff.rfl}
(\lean{EM/Obstruction/RuleTransition.lean}{maxR\_iff},
\lean{EM/Obstruction/RuleTransition.lean}{minR\_iff}).
This is a formal statement separating the two Euclid--Mullin sequences at
the level of congruence obstructions: the class is inhabited under
$\mathrm{maxFac}$ (and machine-checks Cox--van der Poorten's omission of~$5$)
and empty under $\minFac$.  It says nothing about the truth values of the
two conjectures, one of which is open.
\end{keyresult}

\paragraph{Why ``every modulus'' is the load-bearing word.}
An earlier form of the min-side theorem carried a parity hypothesis: the
fragment was proved empty at every \emph{odd} modulus.  That left the
statement open to a fair objection, because the one obstruction anybody
has ever exhibited for a Euclid-type sequence---Cox and van der
Poorten's---lives at $m = 12$.  The technique that works was outside the
theorem that says the technique fails.

The parity entered at a single point.  Extraction has to produce, at each
invariant residue~$r$, an \emph{odd} candidate $N \geq 3$ in the class
$r+1$; for odd~$m$ every class contains odd naturals, but for even~$m$ a
class can consist entirely of even ones, and there the fragment's step and
avoid clauses say nothing.  The repair is to carry the parity as a witness:
$\mathrm{OddRepresentable}(m,r)$ asserts that the class $r+1$ contains an
odd natural.  It holds along the orbit, since $\Prod{2}(n)$ is even; it is
closed under the transition, and for a reason cheaper than expected, since
$r+1$ odd forces $r$ even and $r \cdot s$ stays even whatever the multiplier
is; and it is exactly what realization needs, because $N = c + 2mK$ remains
in the class, remains odd, and exceeds any prescribed bound at
\emph{every} modulus
(\lean{EM/Obstruction/RuleTransition.lean}{exists\_large\_odd\_of\_representable}).
Intersecting the invariant with it yields a certificate at every modulus, and
the oddness hypothesis drops out
(\lean{EM/Obstruction/RuleTransition.lean}{no\_congruence\_induction\_proof\_of\_ne\_zero}).

\paragraph{A min-side obstruction \emph{is} an induction proof.}
The unification also collapses a distinction the two files had maintained
separately.  Under $\Phi = \minFac$, propagation is the step case of an
induction, tail-containment supplies its base, and blocking is precisely its
avoid clause: a min-side obstruction and a congruence-invariant induction
proof are the same object
(\lean{EM/Obstruction/RuleTransition.lean}{toCongruenceProof}).  The semantic
no-go and the proof-theoretic no-go of \S\ref{sec:proof-fragment} are two
readings of one theorem.

\smallskip
\begin{center}
\begin{tabular}{@{}l@{\quad}l@{\quad}l@{}}
\toprule
 & $\boldsymbol{\minFac}$ \textbf{(first)} & $\mathbf{maxFac}$ \textbf{(second)} \\
\midrule
Appearance condition for~$q$
  & open congruence & $q$-smooth, then $q^c$ \\
Hitting set of a missing prime
  & finite & may be infinite \\
Guardian bound $\sigma(N) < q$
  & holds & fails \\
$\sigma(\pi M)$, $\pi < M$ primes
  & $\pi$ & $M$ \\
Reachable set from a free state
  & all units & all units (reorder $\pi, M$) \\
Discarded cofactor
  & large part, unconstrained & small part, already omitted \\
\midrule
\textbf{Congruence obstruction}
  & \textbf{provably none} & \textbf{exists} \\
\bottomrule
\end{tabular}
\end{center}

\smallskip\noindent
The two summary statements are packaged as single theorems:
\lean{EM/Obstruction/NoInvariant.lean}{no\_invariant\_landscape}
and
\lean{EM/Obstruction/MaxVariant.lean}{max\_variant\_landscape}.

\begin{remark}[What this does and does not say about MC]
\label{rem:scope}
The theorem is a no-go, and its scope must be stated precisely.

What is killed is the class of \textbf{elementary congruence
obstructions}: a fixed set of residues modulo a fixed modulus, closed
under the dynamics, containing the orbit tail.  This is exactly the
mechanism that kills primes $\equiv 1 \pmod 4$ on the max side,
including the~$5$ formalized above.  So: \emph{the elementary congruence
mechanism, at every modulus, odd or even, provably cannot exist for the
first sequence.}

What is \textbf{not} killed is Booker's general
argument~\cite{Booker2012}.  It is not a fixed propagating invariant at
all: it is a finite-window reductio.  One assumes that all primes
$\leq X$ except a known exclusion set appear, isolates the last to
appear, and then \emph{constructs} by CRT an auxiliary integer~$d$ whose
Jacobi symbol against the Euclid number is forced to two contradictory
values by quadratic reciprocity.  The construction needs an analytic
existence step at each escalation---Burgess character-sum bounds in
Booker's version, bounds on runs of quadratic residues in the elementary
version of Pollack and Treviño~\cite{PollackTrevino2014}.  The class
handled by Theorem~\ref{thm:no-invariant} does not literally contain
this mechanism.  See \S\ref{sec:reciprocity-frontier}.

Finally, a no-go is not evidence \emph{for} MC beyond its literal
content.  It says only that a disproof of MC cannot be elementary and
congruential.
\end{remark}

%% =========================================================================
\subsection{$\mathrm{IC}_{\min}$ and Its Position in the Network}
\label{sec:ic-min}
%% =========================================================================

It is natural to ask what a congruence-certificate approach to MC would
require.  The formalization states it exactly.

\begin{definition}[$\mathrm{IC}_{\min}$ ---
\lean{EM/Obstruction/NoInvariant.lean}{IC\_min}]
\label{def:ic-min}
\emph{If a prime~$q$ never occurs in the first Euclid--Mullin sequence,
then its omission is certified by a propagating congruence invariant}:
there exist $m \neq 0$ with $\mathrm{RichEnough}(q,m)$ and a set
$S \subseteq \ZZ/m\ZZ$ that is a CvdP obstruction for~$q$.
\end{definition}

\noindent
This is the exact analogue of what is \emph{true} on the max side.  By
Theorem~\ref{thm:no-invariant} it implies MC
(\lean{EM/Obstruction/NoInvariant.lean}{ic\_min\_implies\_mullin}).

\begin{keyresult}
\begin{theorem}[Position in the network ---
\lean{EM/Obstruction/NoInvariant.lean}{ic\_min\_network}]
\label{thm:ic-min-network}
$\DH \Rightarrow \mathrm{IC}_{\min}$,
$\mathrm{SingleHitHypothesis} \Rightarrow \mathrm{IC}_{\min}$, and
\[
  \mathrm{IC}_{\min} \iff \MC .
\]
\end{theorem}
\end{keyresult}

\noindent
\textbf{This must be read correctly.}  The forward implication is the
No-Invariant Theorem; the backward implication
(\lean{EM/Obstruction/NoInvariant.lean}{mullin\_implies\_ic\_min}) is that
MC makes $\mathrm{IC}_{\min}$ \emph{vacuously} true, there being no
missing primes to certify.  Hence $\mathrm{IC}_{\min}$ sits at the
\textbf{top} of the reduction network, strictly above
$\DH$ and $\mathrm{SingleHitHypothesis}$
(\lean{EM/Obstruction/NoInvariant.lean}{dynamicalHitting\_implies\_ic\_min},
\lean{EM/Obstruction/NoInvariant.lean}{singleHit\_implies\_ic\_min}).

$\mathrm{IC}_{\min}$ is therefore \textbf{not} a weakening of MC and
\textbf{not} a new route to a proof.  It is the precise statement of
what a congruence-certificate approach would require, and it is
logically identical to the conjecture itself.  The entire scientific
value of \code{EM/Obstruction/NoInvariant.lean} lies in the no-go content
of \code{no\_cvdp\_obstruction}, not in the reduction $\mathrm{IC}_{\min}
\Rightarrow \MC$.

%% =========================================================================
\subsection{The Obstruction Calculus}
\label{sec:obstruction-calculus}
%% =========================================================================

The proof of Theorem~\ref{thm:no-invariant} barely used the fact that a
certificate tracks \emph{residues}.  Eviction is a statement about which
primes have appeared; fullness is a statement about the richness of the
transition relation; CRT-reach is a statement about forcing states.
Only the last mentions congruences at all.  The natural object is
therefore one level up.

\begin{definition}[Enrichment ---
\lean{EM/Obstruction/Calculus.lean}{Enrichment}]
\label{def:enrichment}
An \defname{enrichment} $E$ is any state space a would-be avoidance
certificate is allowed to track: a type $\mathrm{State}$, an
\defname{observation} map $\mathrm{observe} \colon \NN \to
\mathrm{State}$ giving the enriched shadow of the orbit, a
\defname{transition relation} $\mathrm{Trans}$ under which certificates
must be closed, a proof that the actual orbit steps are transitions, and
a predicate $\mathrm{Forcing}\ q$ marking the states at which capture
of~$q$ is at issue.  The transition relation is deliberately allowed to
be \emph{larger} than the true dynamics---the skeptic's-favour rule of
Definition~\ref{def:min-transition}---which makes non-existence of
certificates a stronger statement.
\end{definition}

\begin{definition}[Certificate, killability ---
\lean{EM/Obstruction/Calculus.lean}{Certificate},
\lean{EM/Obstruction/Calculus.lean}{Killable}]
\label{def:certificate}
A \defname{certificate} for~$q$ in~$E$ is a set of states that is
propagating, contains the orbit tail, and contains no forcing state
for~$q$.  The enrichment is \defname{killable at~$q$} if from every late
orbit state a forcing state is reachable inside $\mathrm{Trans}$.
\end{definition}

Eviction, fullness and reach are exactly a proof method for
killability, and once killability is isolated the no-go is three lines.

\begin{theorem}[Emptiness Theorem ---
\lean{EM/Obstruction/Calculus.lean}{no\_certificate}]
\label{thm:emptiness}
A killable enrichment admits no certificate.
\end{theorem}

\begin{proof}
The tail meets the certificate; from that state a forcing state is
reachable; propagation drags the certificate onto it; this contradicts
blocking.
\end{proof}

The abstraction passes the obvious test: instantiating at the congruence
enrichment (states $\ZZ/m\ZZ$, transition $\mathrm{Transition}\ m$,
forcing $\mathrm{ForcingState}\ q\ m$) and proving \emph{that} enrichment
killable
(\lean{EM/Obstruction/Calculus.lean}{congruence\_killable})
recovers Theorem~\ref{thm:no-invariant} with no further argument
(\lean{EM/Obstruction/Calculus.lean}{no\_cvdp\_obstruction'}).

\paragraph{The framework is not vacuous.}
The max side fits the same shape, and there the Cox--van der Poorten
obstruction \emph{is} a certificate
(\lean{EM/Obstruction/Calculus.lean}{maxCertificate\_five}).  By
Theorem~\ref{thm:emptiness} the max enrichment is therefore \emph{not}
killable
(\lean{EM/Obstruction/Calculus.lean}{max\_not\_killable}).  This is
worth emphasising: ``fullness is min-specific'' is not asserted, it is
\emph{derived}---the very theorem provable for $\mathrm{Transition}\ m$
is refutable for its max analogue, and the room between them is exactly
where an omission proof can live.

\paragraph{Comparison of enrichments.}
Enrichments form a category of sorts: a \defname{refinement}
(\lean{EM/Obstruction/Calculus.lean}{Refines}) projects a finer state
space onto a coarser one, certificates pull back along it
(\lean{EM/Obstruction/Calculus.lean}{Certificate.pullback}) and killability pushes
forward
(\lean{EM/Obstruction/Calculus.lean}{Killable.of\_refines}).  Growing the killed
class is therefore a matter of exhibiting new killable enrichments, and
each one removes one more possible shape of a disproof of~$\MC$.

\paragraph{Time-dependent invariants.}
A certificate fixes one invariant set for all time.  That is a real
restriction, and it is removed by grading: a
\defname{graded certificate}
(\lean{EM/Obstruction/Calculus.lean}{GradedCertificate}) carries a
family $S_n$ of state sets, propagating from stage~$n$ to stage~$n{+}1$
and blocking at every stage.  Ordinary certificates are the constant
families
(\lean{EM/Obstruction/Calculus.lean}{Certificate.toGraded}), so this is
strictly wider.

The emptiness argument survives, but only against a \emph{quantitative}
form of killability: one must know how many transitions the reach takes,
because a $k$-step reach from stage~$n$ lands at stage~$n+k$.  The
congruence enrichment satisfies the counted form in the strongest way---
its witness is a \emph{single} transition, multiplication by the
CRT-chosen unit
(\lean{EM/Obstruction/Calculus.lean}{congruence\_killableIn})---so
nothing is lost, and time-dependent congruence invariants do not block a
missing prime either
(\lean{EM/Obstruction/Calculus.lean}{no\_graded\_certificate\_congruence}).

\begin{remark}[Honesty about completeness]
\label{rem:trace-complete}
Over any killable family, the statement ``every missing prime admits a
certificate somewhere in the family'' is \emph{equivalent} to~$\MC$
(\lean{EM/Obstruction/Calculus.lean}{traceComplete\_iff\_mullin}),
exactly as $\mathrm{IC}_{\min} \Leftrightarrow \MC$; for the congruence
family the two statements literally coincide
(\lean{EM/Obstruction/Calculus.lean}{traceComplete\_congruence\_iff\_ic\_min}).
The calculus therefore produces no new \emph{weaker} route to~$\MC$.
Its value is entirely in the growing no-go content, and the whole
landscape is packaged as
\lean{EM/Obstruction/Calculus.lean}{obstruction\_calculus\_landscape}.
\end{remark}

%% =========================================================================
\subsection{The Proof-Theoretic Dichotomy}
\label{sec:proof-fragment}
%% =========================================================================

Everything so far concerns \emph{semantic} invariants: sets of states.
One can go further and ask about \emph{proofs}.  Fix a proof shape---the
one every classical omission argument in this family has---and ask
whether the shape can establish anything at all on the min side.

\begin{definition}[The congruence-induction fragment ---
\lean{EM/Obstruction/Fragment.lean}{CongruenceInductionProof}]
\label{def:fragment}
A \defname{congruence-invariant induction proof} that~$q$ is eventually
never captured consists of an invariant $\mathrm{inv} \subseteq
\ZZ/m\ZZ$, a stage~$N_0$ at which the actual orbit satisfies it, an
induction \emph{step}, and a \emph{conclusion}.  The step and the
conclusion are required to hold for \emph{every} admissible candidate
$N$ (odd, $\geq 3$) in the class $r+1$: the proof's only knowledge of
the orbit is the residue~$r$, so it cannot distinguish candidates the
invariant does not separate.
\end{definition}

The fragment is a genuine proof system: an inhabitant really does prove
the omission
(\lean{EM/Obstruction/Fragment.lean}{eventually\_avoids}).  It is also
not a straw man in the other direction---its conclusion field asks only
that \emph{capture} be excluded ($\minFac N \neq q$), the weakest
conclusion an omission proof can have, which makes its emptiness a
stronger theorem.

Proof mining converts an inhabitant into a certificate
(\lean{EM/Obstruction/Fragment.lean}{toCertificate}): the step case,
read semantically, is propagation; the conclusion, read semantically,
is blocking.  With Theorem~\ref{thm:emptiness} this gives:

\begin{theorem}[Unprovability ---
\lean{EM/Obstruction/Fragment.lean}{no\_congruence\_induction\_proof}]
\label{thm:unprovability}
For a missing prime~$q$ the fragment is \textbf{empty}, at every odd
modulus.  No richness hypothesis is needed: a proof at a poor modulus
lifts to a rich one
(\lean{EM/Obstruction/Fragment.lean}{lift}), and at the proof level the
$\mathrm{BlocksDeath}$ caveat of the lifting discussion in
\S\ref{sec:no-invariant} dissolves,
because every field of a fragment proof is universally quantified over
candidates and candidates cast down.
\end{theorem}

The converse direction is where the statement acquires its sting.  For a
prime that \emph{has} appeared the fragment is inhabited, and by
Euclid's own argument: once $q \mid \Prod{2}(n)$ it divides the
accumulator forever, so it can never divide a candidate---an induction
whose invariant is the congruence condition ``$q \mid r$''
(\lean{EM/Obstruction/Fragment.lean}{appearedProof}).

\begin{keyresult}
\textbf{Provability decides appearance}
(\lean{EM/Obstruction/Fragment.lean}{congruence\_provability\_iff}).
For a prime~$q$ and an odd modulus~$m$ divisible by~$q$,
\[
  \text{the fragment proves ``$q$ eventually avoided''}
  \quad\Longleftrightarrow\quad
  \text{$q$ appears in the sequence.}
\]
Within this proof genre, provability of avoidance and membership in the
sequence are the \emph{same predicate}: the fragment can certify an
avoidance only for the trivial reason that there is nothing to avoid.
\end{keyresult}

\paragraph{The max-side control.}
The corresponding max fragment is \emph{inhabited} at $q = 5$,
$m = 12$: Cox--van der Poorten's proof is literally a fragment
inhabitant, with invariant $\{6\}$ and base stage~$1$
(\lean{EM/Obstruction/Fragment.lean}{maxProofFive}), and fragment
soundness re-derives the omission
(\lean{EM/Obstruction/Fragment.lean}{max\_fragment\_proves\_five}).  So
the provability equivalence \emph{fails} for the max rule
(\lean{EM/Obstruction/Fragment.lean}{max\_provability\_not\_iff}).  The
dichotomy, now at the level of proofs rather than invariants: \emph{the
same proof system that proves the max sequence misses~$5$ provably
cannot establish any omission for the min sequence}
(\lean{EM/Obstruction/Fragment.lean}{proof\_theoretic\_dichotomy}).

\paragraph{How wide is the fragment?}
Definition~\ref{def:fragment} restricts the invariant to be uniform
in~$n$ and the clauses to be uniform over candidates.  Three relaxations
suggest themselves, and all three are free.

\begin{itemize}[leftmargin=*]
\item \textbf{Time-dependence.}  Let the invariant be a family
  $\mathrm{inv}_n$.  This dies because the congruence enrichment is
  killable in one step (\S\ref{sec:obstruction-calculus}), so a graded
  invariant is dragged onto a forcing state exactly as a constant one
  is.
\item \textbf{Size.}  Let the step and conclusion be required only for
  candidates $N \geq B(n)$, so the proof may use an archimedean lower
  bound on the Euclid number.  This dies because Dirichlet hands out
  arbitrarily large primes: the two-prime construction of
  Move~2 can be run with its first prime above any prescribed bound
  (\lean{EM/Obstruction/NoInvariant.lean}{free\_transition\_large}), and
  a class always contains arbitrarily large odd representatives
  (\lean{EM/Obstruction/NoInvariant.lean}{exists\_large\_odd\_in\_class}).
\item \textbf{Factor count.}  Let the clauses be required only for
  candidates with at least $K(n)$ distinct prime factors.  This dies for
  a reason worth isolating: multiplying by a prime $p \equiv 1
  \pmod m$ chosen above the current value changes \emph{neither the
  residue class nor the least prime factor}, while raising $\omega$
  by one
  (\lean{EM/Obstruction/NoInvariant.lean}{exists\_class\_omega}).
  Iterating pushes $\omega$ as high as desired without disturbing
  anything the congruence machinery can see.
\end{itemize}

A guard is \defname{admissible} when it admits the orbit's own
candidates---$B(n) \leq \Prod{2}(n) + 1$ and $K(n) \leq
\omega(\Prod{2}(n)+1)$---since otherwise the fragment says nothing about
the orbit and soundness fails.  Note how generous that is: the proof may
assume the candidate is as large as the Euclid number itself, and has as
many distinct prime factors as it actually has.

\begin{theorem}[The widened unprovability theorem ---
\lean{EM/Obstruction/Fragment.lean}{no\_omega\_graded\_induction\_proof}]
\label{thm:widened-unprovability}
For a missing prime~$q$ and any admissible guards $(B, K)$, the fragment
with a stage-dependent invariant and both guards is still empty at every
odd modulus; and provability still decides appearance
(\lean{EM/Obstruction/Fragment.lean}{omega\_graded\_provability\_iff}).
The eight clauses are packaged as
\lean{EM/Obstruction/Fragment.lean}{graded\_proof\_theoretic\_dichotomy}.
\end{theorem}

\paragraph{The smoothness axis, and why it closes too.}
It is tempting to record ``multiplicative anatomy'' wholesale as the
surviving guard direction, on the grounds that the fullness construction
produces $N = \pi M$ with $\omega(N) = 2$.  The third bullet above
refutes that: $\omega$ can be raised arbitrarily at no cost.

The genuinely different direction is the \emph{opposite} one---a guard
bounding the candidate from above, demanding it be $y$-smooth.  Every
construction in this section produces candidates whose cofactors are
huge primes, so such a guard is never \emph{reached} by the killing
argument.  But it is also never \emph{usable}, and for an independent
reason: a guard is meaningful only if the orbit's own candidates satisfy
it, and they do not.  Each prime either enters the accumulator or has a
finite hitting set, so for every~$y$ the Euclid numbers are eventually
$y$-rough---every prime factor exceeds~$y$
(\lean{EM/Obstruction/NoInvariant.lean}{eventually\_rough}).  A
$y$-smooth fragment therefore excludes the orbit's own candidates from
some stage on, and proves nothing about the orbit
(\lean{EM/Obstruction/NoInvariant.lean}{smooth\_guard\_inadmissible},
\lean{EM/Obstruction/Fragment.lean}{smooth\_fragment\_never\_sound}).

\begin{keyresult}
\textbf{The fragment analysis is complete for fixed-threshold guards}
(\lean{EM/Obstruction/Fragment.lean}{fragment\_analysis\_complete}).
Guards bounding the candidate from \emph{below}---stage-dependence,
size, number of prime factors---are free: the fragment carrying all of
them is still empty for every missing prime, and provability still
decides appearance.  The guard bounding it from \emph{above} is
unusable: no fixed smoothness threshold is admissible.
\end{keyresult}

\begin{frontierbox}
\textbf{The growing guard closes too.}  A \emph{growing} smoothness
guard $y(n)$ looks like a genuine residue: its admissibility is the
anatomy statement $P^{+}(\Prod{2}(n)+1) \leq y(n)$, and the
constructions above produce candidates with large prime factors.  It is
not, and the reason is an order-of-quantifiers point.

The killing construction does not need candidates tailored to the stage.
By pigeonhole the orbit returns to some residue $r_0$ modulo the working
modulus cofinally often
(\lean{EM/Obstruction/Fragment.lean}{exists\_recurrent\_residue}), and
the candidates' defining conditions mention the stage \emph{only}
through that residue.  So one may choose the candidates first and the
stage afterwards.  Once fixed they are two specific integers with some
largest prime factor~$C$, and admissibility together with
\code{eventually\_rough} forces $y(n) > C$ at all late stages.  Every
admissible growing guard therefore admits them
(\lean{EM/Obstruction/Fragment.lean}{no\_smooth\_graded\_induction\_proof}).

\smallskip
\noindent
The guard axis is therefore closed in both directions
(\lean{EM/Obstruction/Fragment.lean}{guard\_analysis\_complete}): weakening
the proof's obligations---in the size of the candidate, its number of
prime factors, or its smoothness---never helps.

\smallskip\noindent
Stated positively, and within the fragment analysed here: \emph{any
omission proof for the $\minFac$ sequence whose knowledge of the
candidate is congruential---fixed or growing modulus, with the guards
above---must first establish an unproven anatomy property of the Euclid
numbers}, specifically about the large prime factors of
$\Prod{2}(n)+1$, exactly the part that $\minFac$ discards.  Proofs
outside this fragment are not excluded by anything here.
\end{frontierbox}

%% =========================================================================
\subsection{The Reciprocity Frontier}
\label{sec:reciprocity-frontier}
%% =========================================================================

Remark~\ref{rem:scope} left a gap: Booker's mechanism is not a fixed
congruence invariant, so Theorem~\ref{thm:no-invariant} does not
literally cover it.  That gap is now closed \emph{for induction proofs
with a $\Pi_n$-congruence-closed invariant} (Booker's and
Pollack--Trevi\~no's actual arguments are finite-window reductios using
smoothness of the largest-factor sequence, and nothing shows a min-side
transplant would have the invariant shape).  The symbol algebra was
machine-verified first---the Euclid unit and parity laws, the
moving-modulus symbol update, and the Character Non-Constancy Lemma---and
the assembled no-invariant statement is now a theorem as well
(Theorem~\ref{thm:no-reciprocity-invariant}).

\paragraph{Reciprocity invariants are congruence invariants with a
moving modulus.}
For a \emph{fixed} finite anchor set $d_1,\dots,d_s$, each Jacobi symbol
$\left(\tfrac{d_j}{\,\cdot\,}\right)$ is a real Dirichlet character, so
tracking it is tracking the accumulator modulo a fixed modulus:
these are already covered by Theorem~\ref{thm:no-invariant} at
$m' = \mathrm{lcm}(m, 4|d_1|, \dots, 4|d_s|)$.  The genuinely new data
are symbols whose \emph{modulus is the new multiplier}: for orbit primes
$p_i$, the bits $\left(\tfrac{p_i}{p_j}\right)$ together with
$p_i \bmod 8$.  The whole symbol update at step~$n$ is a function of the
new multiplier modulo
\[
  \Pi_n \;=\; 8\, m\, \Prod{2}(n),
\]
and this is a theorem, not a heuristic:
\lean{EM/Reciprocity/SymbolAlgebra.lean}{symbolModulus\_spec} shows that
two odd candidates agreeing mod~$\Pi_n$ carry identical level-$n$
reciprocity data---every symbol $J(p_i \mid \cdot)$ for $i \leq n$, the
class mod~$8$, and the class mod~$m$.  So the reciprocity enrichment is
\emph{exactly} the replacement of a fixed modulus by an orbit-dependent
growing one, with old coordinates never rewritten.

This is what makes the formalization tractable: a reciprocity invariant
need not be modelled as a tuple of symbol coordinates with a
realisability side condition.  It is a predicate on accumulators closed
under congruence modulo~$\Pi_n$
(\lean{EM/Reciprocity/NoReciprocityInvariant.lean}{ReciprocityInductionProof}),
which is a statement of first-order shape with no dependent types.

\paragraph{Why the three moves survive---concretely.}
\emph{Eviction is automatic.}  $\Pi_n$ contains $\Prod{2}(n)$ as a
factor, and the Euclid number is coprime to it; $\Pi_n$ contains~$8$,
and the Euclid number is odd.  So the Euclid unit law
($N \equiv 1 \bmod p_i$) and the mod-$4$ law ($N \equiv 3 \bmod 4$) are
\emph{consequences} of taking the candidate in the class
$\Prod{2}(n)+1$ mod~$\Pi_n$, not additional hypotheses.  The orbit
primes evict themselves.

\emph{Fullness needs no new analytic input.}  Free-state fullness
(Move~2) was proved for an \emph{arbitrary} modulus, so it applies
verbatim at~$\Pi_n$
(\lean{EM/Obstruction/NoInvariant.lean}{free\_transition\_large}):
Dirichlet at~$\Pi_n$ supplies the new multiplier in any prescribed
class, and the cofactor realises the candidate's class.  The
qualitative Dirichlet input is the same theorem at a larger modulus---no
Chebotarev, no Burgess.

\emph{Reach is unchanged.}  Forcing is a condition modulo~$m$ alone and
$m \mid \Pi_n$, so the forcing states of the fixed-modulus theory are
still forcing at every stage.  The one piece of bookkeeping is that the
CRT-chosen unit lives mod~$m$ and must be lifted to a unit mod~$\Pi_n$,
which is possible because $\mathrm{unitsMap}$ is surjective along
divisibility.

\begin{theorem}[Reciprocity-class No-Invariant Theorem ---
\lean{EM/Reciprocity/NoReciprocityInvariant.lean}{no\_reciprocity\_induction\_proof}]
\label{thm:no-reciprocity-invariant}
Let $q$ be a prime that never occurs in the first Euclid--Mullin
sequence and let $m \neq 0$.  Then there is \textbf{no} induction proof
of $q$'s eventual avoidance whose invariant is closed under congruence
modulo the growing symbol modulus $\Pi_n = 8\,m\,\Prod{2}(n)$.
Consequently, in this fragment as in the fixed-modulus one, provability
of avoidance \emph{decides} appearance
(\lean{EM/Reciprocity/NoReciprocityInvariant.lean}{reciprocity\_provability\_iff}):
the fragment certifies an avoidance only when there is nothing to avoid.
\end{theorem}

\paragraph{What the Character Non-Constancy Lemma contributes.}
The formal proof above runs through fullness at~$\Pi_n$ and does not
invoke character non-constancy.  The lemma remains the conceptual
account of \emph{why} the max-side argument cannot transfer, and it too
is machine-verified: a nontrivial character on $\ZZ/D\ZZ$ takes distinct
values on the primes above any bound~$Y$
(\lean{EM/Reciprocity/SymbolAlgebra.lean}{char\_non\_constancy},
with the eventually-constant variant
\lean{EM/Reciprocity/SymbolAlgebra.lean}{char\_not\_eventually\_constant}).
Under $\mathrm{maxFac}$ the factor support of the Euclid number is
confined to a \textbf{finite} set of primes, where a character can be
made constant---this is exactly what Booker's and Pollack--Treviño's
choice of the auxiliary discriminant~$d$ achieves, and the bottom-up
product then contradicts the top-down evaluation.  Under $\minFac$ the
support is confined to a \textbf{cofinite} set, where no character is
ever constant, so no choice of~$d$ produces a contradiction.  Two proofs
of the same asymmetry: one structural, one formal.

\begin{frontierbox}
\textbf{The payoff, and its exact scope.}  Two axes of widening have to
be distinguished, and the results close one of them completely.

\smallskip
\noindent\emph{Obligations.}  Which candidates the step and conclusion
clauses must handle.  Weakening these---by a lower bound on the size of
the candidate, a lower bound on its number of prime factors, or an upper
bound on its prime factors---never helps
(\S\ref{sec:proof-fragment},
\lean{EM/Obstruction/Fragment.lean}{guard\_analysis\_complete}).
\textbf{Closed.}

\smallskip
\noindent\emph{State.}  What the invariant may track.  Residues modulo a
fixed modulus (Theorem~\ref{thm:no-invariant}) and modulo the growing
symbol modulus (Theorem~\ref{thm:no-reciprocity-invariant}) are both
killed.  \textbf{Closed for congruence data.}

\smallskip\noindent
What remains is a state tracking \emph{anatomy} rather than congruence
data---$\omega(\Prod{2}(n)+1)$, its largest prime factor, its
smoothness---which is a different object from a smoothness \emph{guard}
and is not covered by any theorem here.  That is where the max-side
mechanism lives: Cox--van der Poorten's non-congruence ingredient is the
$q$-smoothness forced by $\mathrm{maxFac}$, and the live min-side
example is the eventually-prime branch of
\S\ref{sec:autonomous-branch}.  It is also, so far as we know, where any
disproof of MC would have to live.
\end{frontierbox}

\noindent
A second boundary is more radical still and lies outside every framework
in this section: \emph{exact-value} arguments, which use
$\Prod{2}(n)+1$ as an integer identity rather than as a congruence
class.

\smallskip\noindent
Two expectations recorded in earlier drafts turned out to be too
pessimistic, and it is worth saying so.  That $\omega$ would survive as
an axis: it does not---raising $\omega$ costs the killing argument
nothing.  That a growing smoothness guard would survive: it does not
either---the candidates can be chosen before the stage.  What survives
is neither of those, but the strictly different possibility of an
invariant whose \emph{state} is anatomical rather than congruential.
That possibility is settled in the next subsection.

%% =========================================================================
\subsection{The Anatomy Axis, Settled}
\label{sec:anatomy-axis}
%% =========================================================================

An omission proof reasons backwards from the selection event: \emph{if
$q$ were selected at $N$, then $N$ would have to look like this}.  So the
content of the axis is the strength of the implication
\[
  \Phi(N) = q \;\Longrightarrow\; \text{(some property of } N\text{)}.
\]
For the two rules the answers are opposite, and the gap between them is
exactly the gap between congruence data and anatomy.

\begin{theorem}[$\minFac$ selection is a congruence condition ---
\lean{EM/Obstruction/Anatomy.lean}{minFac\_eq\_congruence\_determined}]
\label{thm:min-congruence-determined}
For an odd $N \geq 3$ and a prime $q$,
\[
  \minFac(N) = q
  \iff
  q \mid N \ \text{ and } \ p \nmid N \text{ for every odd prime } p < q .
\]
Every clause is a divisibility by a prime $\leq q$.  Consequently, if $M$
is rich enough for $q$ and $N_1 \equiv N_2 \pmod M$ are both odd and
$\geq 3$, then $\minFac(N_1) = q$ if and only if $\minFac(N_2) = q$.
\end{theorem}

\begin{theorem}[$\mathrm{maxFac}$ selection is determined by no modulus ---
\lean{EM/Obstruction/Anatomy.lean}{maxFac\_not\_congruence\_determined}]
\label{thm:max-not-determined}
For every modulus $M \neq 0$ there are odd $N_1, N_2 \geq 3$ with
$N_1 \equiv N_2 \pmod M$ and a prime $q$ such that
$\mathrm{maxFac}(N_2) = q$ and $\mathrm{maxFac}(N_1) \neq q$.
\end{theorem}

\noindent
The witness is cheap, which is itself the point.  Take $N_1 = 2M+1$ and,
by Dirichlet, a prime $p \equiv 1 \pmod M$ with $p > N_1$.  Then
$\mathrm{maxFac}(p) = p$ while $\mathrm{maxFac}(N_1) \leq N_1 < p$.  No
amount of congruence bookkeeping sees the difference.

So the min rule hands an omission proof \emph{no anatomy at all}: its
selection condition is already inside the congruence fragment, which
\S\ref{sec:formal-dichotomy} shows to be empty at every modulus.  The max
rule hands over genuine anatomy---$\mathrm{maxFac}(N) = 5$ forces
$5$-smoothness---and that is the ingredient Cox--van der Poorten's
argument consumes at its second move.

\paragraph{Anatomy carried as state.}
There remains the possibility of an invariant whose anatomy component is
not derived from the selection event at all: $\omega$ of the accumulator,
its largest prime factor, its smoothness.  Model it directly.  An
\emph{anatomy-state induction proof} has an invariant on pairs
(residue, anatomy value), and its step clause must admit \emph{every}
anatomy value---because the anatomy of $\Prod{2}(n+1)$ is not a function
of the residue of $\Prod{2}(n)$ together with the candidate.  A proof
that could predict it would already be an anatomy theorem about the
Euclid numbers, which is the thing being sought, not a tool for finding
it.

\begin{theorem}[Anatomy as state is inert ---
\lean{EM/Obstruction/Anatomy.lean}{no\_anatomy\_induction\_proof}]
\label{thm:anatomy-inert}
For every missing prime $q$, every modulus $m \neq 0$ and every anatomy
type, the anatomy-state fragment is empty.
\end{theorem}

The proof is a projection: forgetting the anatomy component leaves a
congruence-invariant induction proof
(\lean{EM/Obstruction/Anatomy.lean}{AnatomyInductionProof.toCongruenceProof}).
The step survives because the anatomy value may be carried unchanged, and
the avoid clause never consulted it---whether $q$ is selected at $N$
depends on $N$, not on how the accumulator factors.

\begin{frontierbox}
\textbf{What the anatomy axis leaves.}
Nothing on this axis is an \emph{invariant} question any more.  What
survives is a demand for a theorem about the anatomy of the specific
integers $\Prod{2}(n)+1$: that they are composite infinitely often
(\S\ref{sec:composite-floor}), that their largest prime factor is large,
that they are not smooth.  Those are not properties of a proof system;
they are number-theoretic facts about one orbit, and the
orbit-specificity barrier applies to them verbatim.

The residue of the obstruction programme is therefore not a wider class
of invariants to kill.  It is a single, named, open anatomy statement.
\end{frontierbox}
